Desastres naturais provocados por eventos climáticos extremos causam impactos profundos e severos de diversas maneiras.
No aspecto econômico, os recursos alocados para a recuperação de áreas afetadas enfraquecem a microeconomia local,
tornando a região ainda mais vulnerável a eventos futuros. Estruturalmente, o tempo e o dinheiro investidos na
construção e no aprimoramento de edificações, tanto privadas (como residências e negócios) quanto públicas (como
infraestrutura urbana e áreas administrativas), podem ser comprometidos pela falta de preparação para lidar com tais
ocorrências. Já o impacto social é, sem dúvida, o mais significativo. Este impacto, imensurável, é sentido por diversas
gerações de pessoas afetadas por desastres, como evidenciado pelos efeitos do furacão Katrina\cite{reardon}\cite{neria}.

Devido às mudanças climáticas, a frequência e intensidade dos desastres naturais vêm aumentando gradualmente desde os
anos 1950\cite{dunn}. Como resultado, muitas regiões estão sendo afetadas por eventos climáticos extremos de maneira
cada vez mais rápida e despreparada. Entre esses desastres, as enchentes se destacam pela sua capacidade de causar
destruição generalizada e afetar áreas densamente povoadas, especialmente em regiões urbanas com infraestrutura
inadequada para gerenciar grandes volumes de água. No final de abril e início de maio de 2024, o estado do Rio Grande do
Sul passou por um período de alta precipitação que resultou na cheia das suas principais bacias hidrográficas. Na Região
Metropolitana de Porto Alegre, esse cenário culminou na maior cheia histórica do Rio Guaíba, que ultrapassou a marca de
4.76 metros, superando a cheia de 1941, até então considerada a mais devastadora da cidade\cite{dep}. Esse desastre foi
impulsionado por uma combinação de fatores meteorológicos, e, combinado com a maior vulnerabilidade de centros urbanos à
enchentes\cite{yang}\cite{hammond}\cite{zhang}\cite{pereira}, levou à um evento de devastação generalizada, com milhares
de pessoas deslocadas e centenas de vítimas fatais\cite{alcantara}.

% Falar das mudanças geográficas em POA?

Durante a gestão da crise, o monitoramento e predição dos níveis de água do Rio Guaíba – o mais importante da região –
foi fundamental para orientar os tomadores de decisão em suas ações de planejamento e para que à população pudesse
reagir melhor à ameaça iminente. Essas previsões foram feitas principalmente através de modelos hidrológicos e
hidrodinâmicos, que, embora precisos, são complexos, exigem grande quantidade de dados topográficos, meteorológicos e
hidrológicos, requerem conhecimento especializado em hidrologia e demandam tempo de simulação prolongado. Esse último
fator foi acentuado nos momentos mais delicados da crise, quando as previsões só podiam ser realizadas com uma
frequência diária \cite{iph}, o que as tornou insuficientes para atender às demandas de previsões em tempo real ou de
curto prazo \cite{atashi}.

Uma alternativa que visa superar essas limitações são os modelos baseados em dados. Os mesmos utilizam técnicas para
identificar e explorar relações estatísticas entre os dados de entrada e saída. Com a prendizagem dos padrões presentes
nos dados históricos, é possível mapear as dinâmicas subjacentes do sistema e utilizá-las para prever valores futuros.
Assim, esses modelos são capazes de realizar previsões precisas com base em séries temporais, capturando tendências e
comportamentos complexos que podem ser difíceis de modelar diretamente por meio de abordagens tradicionais. Modelos
baseados em dados são mais rápidos e eficientes para previsões em tempo real, sendo uma escolha adequada para sistemas
de alerta precoce e mitigação de desastres de cheias\cite{ta}.

% Ligar melhor a explicação dos modelos com o problema das enchentes para que o leitor não esqueça do que foi definido;

Uma das maiores vantagens de utilizar essa abordagem é o grande desenvolvimento do campo de Análise e Predição de Séries
Temporais. Com isso, é possível utilizar a grande quantidade de técnicas desenvolvidas para a aplicação em problemas com
uma modelagem semelhante, mas em contextos distintos \cite{hyndman}, como Economia, Biologia, entre outros. Dentro desta
abordagem, uma ampla gama de modelagens são utilizadas para a predição de séries temporais. Isso apresenta uma
oportunidade de análise e comparação entre vantagens e desvantagens das mesmas com o intuíto de descobrir a mais efetiva
para a situação particular \cite{makridakis}\cite{morales}.

Uma abordagem mais clássica e altamente implementada destes modelos vem na na forma de modelos estatísticos, como o
ARIMA (\textit{Autoregressive Integrated Moving Average}) \cite{box}, um modelo que descreve processos estocásticos
usando uma combinação de autoregressão, isto é, a predição de uma Série Temporal usando somente valores anteriores da
mesma, a média móvel - ajuste das previsões considerando os erros residuais das observações anteriores, o que permite
capturar flutuações - e, por último, a diferenciação dos valores consecutivos da série, reduzindo o impacto de padrões
de longo prazo e tendências. Além disso, um elemento de sazionalidade, ou seja, padrões repetidos em com uma certa
frequência fixa, como padrões mensais ou trimestrais, podem ser compreendidos através de uma extensão do ARIMA, o
SARIMA. 

% Se for usar o ETC ou o DOT, aqui é o ponto para mencionar eles;

Esses modelos, eficazes e de baixo custo computacional, apresentam tanto limitações quanto vantagens quando comparados
com modelos mais recentes desenvolvidos nas áreas de Aprendizado de Máquina (\textit{Machine Learning}) e Aprendizado
Profundo (\textit{Deep Learning}) \cite{kontopoulou}. No âmbito de \textit{Machine Learning}, modelos baseados em
\textit{Gradient Boosting} têm se mostrado extremamente eficazes. Um exemplo amplamente utilizado é o XGBoost
(\textit{eXtreme Gradient Boosting}) \cite{chen_xgboost}. O mesmo utiliza árvores de decisão em série, ajustando
constantemente os erros das previsões anteriores. Além disso, ele incorpora regularização para reduzir o risco de
\textit{overfitting} e otimizações específicas para lidar com grandes volumes de dados. Uma implementação semelhante
visando reduzir o tempo computacional é dada pelo LightGBM \cite{lightgbm}. Ela mantém as vantagens do XGBoost, como
otimização para dados esparsos, treinamento paralelo, múltiplas funções de perda e regularização, mas altera o método de
construção das árvores de decisão.

Por outro lado, em Aprendizado Profundo, modelos baseados em Redes Neurais Recorrentes (RNNs), como o \textit{Long Short
Term Memory} (LSTM), são amplamente utilizados em situações semelhantes, como demonstrado por Borwarnginn et al.
\cite{borwarnginn}. O LSTM foi projetado especificamente para lidar com dependências de longo prazo em séries temporais,
garantindo que informações relevantes de eventos passados sejam preservadas, ao mesmo tempo que mantém o contexto de
eventos mais recentes. Essa abordagem é particularmente útil em problemas onde padrões temporais complexos e não
lineares precisam ser modelados.

Além dos modelos baseados em RNNs, outra abordagem amplamente utilizada em Aprendizado Profundo são as baseadas em Redes
Neurais Perceptron Multicamadas (MLP, \textit{Multilayer Perceptron}). Diferentemente das RNNs, que são projetadas para
processar dados sequenciais, as MLPs utilizam camadas conectadas para identificar padrões em conjuntos de dados
estruturados. Um exemplo é o modelo NHITS (\textit{Neural Hierarchical Interpolation for Time Series Forecasting})
\cite{challu}, que adapta a arquitetura MLP para o contexto de séries temporais. O NHITS emprega uma estratégia
hierárquica para interpolar e extrair características relevantes em diferentes escalas temporais, permitindo capturar
tanto padrões de curto quanto de longo prazo. Com isso, o mesmo apresenta uma grande capacidade de lidar com séries
temporais irregulares ou com alta variabilidade.

Com todas essas diferentes abordagens em mente, este estudo visa a aplicação dos mesmos no mesmo contexto de predição do
nível da água nas enchentes do Rio Guaíba em 2024 e a comparação entre os mesmos será feita com base em métricas de
desempenho como o \textit{Root Mean Squared Error} (RMSE), \textit{Mean Absolute Error} (MAE) e o \textit{Nash-Sutcliffe
Model Efficiency Coefficient} (NSE).

Para isso, serão utilizados dados das séries temporais do Sistema Nacional de Informações sobre Recursos Hídricos
(SNIRH), que disponibiliza informações das principais bacias hidrográficas do Brasil\cite{snirh}. Esses dados serão a
base para as séries históricas do nível de água, complementados por dados externos, como os do MERRA-2
(\textit{Modern-Era Retrospective analysis for Research and Applications, Version 2}), um projeto da NASA que integra
observações de satélites e estações terrestres para gerar dados meteorológicos globais\cite{nasa}. Dados meteorológicos
adicionais serão obtidos a partir de fontes públicas, como o Instituto Nacional de Meteorologia (INMET) e o Instituto
Brasileiro de Geografia e Estatística (IBGE)\cite{inmet}\cite{ibge}.